% theme: earthquakes
\MajorQuestion{地震}
\SetRowSpaceQanda

\Instruction{揺れが始まる場所と、その観測に関する次の問いに答えなさい。}
\QQ{地下で岩石の破壊が最初に起こった地点を\NamedBlank{hypocenter}という。}{}
\QQ{\RefBlank{hypocenter}の真上にあたる地表の地点を\NamedBlank{epicenter}という。}{}
\QQ{地面の揺れを記録する装置を何というか。}{}
\QQ{はじめに伝わる小さな揺れを\NamedBlank{initial_tremor}という。}{}
\QQ{\RefBlank{initial_tremor}のあとに伝わる大きな揺れを\NamedBlank{main_tremor}という。}{}

\Instruction{伝わる波と、記録から分かる量に関する次の問いに答えなさい。}
\QQ{観測点に先に到着し、\RefBlank{initial_tremor}を起こす波を\NamedBlank{p_wave}という。}{}
\QQ{観測点にあとから到着し、\RefBlank{main_tremor}を起こす波を\NamedBlank{s_wave}という。}{}
\QQ{\RefBlank{initial_tremor}が始まってから、\RefBlank{main_tremor}が始まるまでの時間を何というか。}{}
\QQ{ある場所における揺れの大きさを表す尺度を何というか。}{}
\QQ{地下で起こった現象の規模を、放出されたエネルギーの大きさで表す値を何というか。}{}

\Instruction{日本付近で起こる現象と、地表をおおう岩盤に関する次の問いに答えなさい。}
\QQ{地表をおおい、少しずつ移動している巨大な岩盤を\NamedBlank{plate}という。}{}
\QQ{海側の\RefBlank{plate}が陸側のものの下へ沈みこむ境界付近で起こる地震を何というか。}{}
\QQ{陸側の\RefBlank{plate}内部にある岩盤のずれが動くことで起こる地震を何というか。}{}
\QQ{過去に繰り返し動き、今後も動く可能性がある、岩盤にできたずれを何というか。}{}
\QQ{海側の\RefBlank{plate}が陸側のものの下へ沈みこむ境界に沿ってできる、細長く深い地形を何というか。}{}

\Instruction{大地の変化による地形や災害と、防災に関する次の問いに答えなさい。}
\QQ{インド半島をのせた\RefBlank{plate}の衝突によって、海底の地層が押し上げられてできた山脈を何というか。}{}
\QQ{海底の急な変化で海水が持ち上げられ、波となって広がる現象を何というか。}{}
\QQ{水を多く含む砂地で、揺れにより水や砂が地表に噴き出す現象を何というか。}{}
\QQ{大地が広い範囲で持ち上がる現象を何というか。}{}
\QQ{\RefBlank{p_wave}を検知し、強い揺れが来る前に知らせる情報を何というか。}{}
